\section{\label{sec:intro}Introduction}

A vital component of such studies is that detection of inverse beta decay (IBD) interactions is difficult due to the very low cross section. %relatively low amount ($\sim$ few MeV) of deposited energy, the reactor signal can be completely swamped by backgrounds.
Almost all reactors are built near the Earth's surface (and near cooling / water resources) for obvious practical reasons.
The optimal local-monitoring detector would be located directly under the reactor it is observing, but such an arrangement is rarely practical.
The observation of the neutrino spectrum leads to the possibility of blind neutrino source range determination (without foreknowledge of the source or distance) \cite{Smith:2010zzb} \cite{Snowmass:2021}
Given adequate statistics (several thousand events), the Fourier transform of the inverse neutrino spectrum (on distance divided by energy) yields a unique measure of the source-detector distance.
To make this practical, one needs both excellent IBD energy resolution (in the range of a few percent) and a low energy threshold (going well below the spectrum peak around 4 MeV, usually background-limited). \cite{Jocher:2013gta}
We do not focus on these matters herein but concentrate on the directionality issue.

The structure of this paper is as follows: 
\S\ref{sec:intro} explains the problem at hand and serves as an introduction.
%\S\ref{sec:dirn} covers how to find directionality in the context of Inverse Beta Decay.
%\S\ref{sec:detectors} covers detector geometries and hybrid setups.
%\S\ref{sec:method} is methodology, including simulations and addressing backgrounds.
%\S\ref{sec:results} presents our results.
%\S\ref{sec:newdesign}  considers new designs.
\S\ref{sec:deltaphi} discusses our novel algorithm for determining direction and resolution.
%In \S\ref{sec:future} we consider future directions of work, and 
We overview and conclude this paper in \S\ref{sec:concl}.
